% GENERATED FILE. Edit canonical lessons, then run npm run generate:books.
% canonical-source-hash: fnv1a64:ba1c8e8e8a3ade34
% canonical-lessons: RU-C215-konechno, RU-C215-navernoye, RU-C215-mozhet-byt, RU-C215-k-sozhaleniyu, RU-C215-k-schastyu

\chapter{Of course, Probably, Maybe, Unfortunately, Fortunately}
\label{ch:ru-of-course-probably-maybe-unfortunately-fortunately}
\hlchaptermodality{215}

\begin{chapteropening}
\textbf{By the end of this chapter:} \emph{I can say of course, probably, maybe, unfortunately and fortunately.}

Complete the last lesson of chapter 215: I can say of course, probably, maybe, unfortunately and fortunately.
\end{chapteropening}

\section[\texorpdfstring{\ru{конечно} (konéchno)}{конечно (konéchno)}]{\ru{конечно} (konéchno) — of course}

\label{lesson:RU-C215-konechno}



\begin{quote}
Before the new one: say the Russian for up, then the Russian for down.
\end{quote}

\subsection*{You'll want to know: \ru{конечно}}
\textbf{\ru{конечно}} (\emph{konéchno}) — \textquotedblleft{}of course\textquotedblright{}.

\begin{grammarlens}[title={What you've built}]
One more word for this chapter.
\end{grammarlens}

\subsection*{Your turn}
\begin{itemize}
\raggedright
  \item \emph{Say it:} \emph{konéchno}
  \item \emph{Say it:} \emph{konéchno}, once more
  \item \emph{Say it:} say \emph{konéchno}
  \item \emph{Recall:} say the Russian for up, then the Russian for down, then say \emph{konéchno} again
\end{itemize}

\begin{tcolorbox}[breakable,colback=teal!4,colframe=teal!35!black,title={Before you move on}]
What does \ru{конечно} mean? (\textquotedblleft{}Of course\textquotedblright{}.) Say it once more.
\end{tcolorbox}
\section[\texorpdfstring{\ru{наверное} (navérnoye)}{наверное (navérnoye)}]{\ru{наверное} (navérnoye) — probably}

\label{lesson:RU-C215-navernoye}



\begin{quote}
Before the new one: say the Russian for down, then the Russian for of course.
\end{quote}

\subsection*{You'll want to know: \ru{наверное}}
\textbf{\ru{наверное}} (\emph{navérnoye}) — \textquotedblleft{}probably\textquotedblright{}.

\begin{grammarlens}[title={What you've built}]
One more word for this chapter.
\end{grammarlens}

\subsection*{Your turn}
\begin{itemize}
\raggedright
  \item \emph{Say it:} \emph{navérnoye}
  \item \emph{Say it:} \emph{navérnoye}, once more
  \item \emph{Say it:} say \emph{navérnoye}
  \item \emph{Recall:} say the Russian for down, then the Russian for of course, then say \emph{navérnoye} again
\end{itemize}

\begin{tcolorbox}[breakable,colback=teal!4,colframe=teal!35!black,title={Before you move on}]
What does \ru{наверное} mean? (\textquotedblleft{}Probably\textquotedblright{}.) Say it once more.
\end{tcolorbox}
\section[\texorpdfstring{\ru{может} \ru{быть} (mózhet byt)}{может быть (mózhet byt)}]{\ru{может} \ru{быть} (mózhet byt) — maybe}

\label{lesson:RU-C215-mozhet-byt}



\begin{quote}
Before the new one: say the Russian for of course, then the Russian for probably.
\end{quote}

\subsection*{You'll want to know: \ru{может} \ru{быть}}
\textbf{\ru{может} \ru{быть}} (\emph{mózhet byt}) — \textquotedblleft{}maybe\textquotedblright{}.

\begin{grammarlens}[title={What you've built}]
One more word for this chapter.
\end{grammarlens}

\subsection*{Your turn}
\begin{itemize}
\raggedright
  \item \emph{Say it:} \emph{mózhet byt}
  \item \emph{Say it:} \emph{mózhet byt}, once more
  \item \emph{Say it:} say \emph{mózhet byt}
  \item \emph{Recall:} say the Russian for of course, then the Russian for probably, then say \emph{mózhet byt} again
\end{itemize}

\begin{tcolorbox}[breakable,colback=teal!4,colframe=teal!35!black,title={Before you move on}]
What does \ru{может} \ru{быть} mean? (\textquotedblleft{}Maybe\textquotedblright{}.) Say it once more.
\end{tcolorbox}
\section[\texorpdfstring{\ru{к} \ru{сожалению} (k sozhaléniyu)}{к сожалению (k sozhaléniyu)}]{\ru{к} \ru{сожалению} (k sozhaléniyu) — unfortunately}

\label{lesson:RU-C215-k-sozhaleniyu}



\begin{quote}
Before the new one: say the Russian for probably, then the Russian for maybe.
\end{quote}

\subsection*{You'll want to know: \ru{к} \ru{сожалению}}
\textbf{\ru{к} \ru{сожалению}} (\emph{k sozhaléniyu}) — \textquotedblleft{}unfortunately\textquotedblright{}.

\begin{grammarlens}[title={What you've built}]
One more word for this chapter.
\end{grammarlens}

\subsection*{Your turn}
\begin{itemize}
\raggedright
  \item \emph{Say it:} \emph{k sozhaléniyu}
  \item \emph{Say it:} \emph{k sozhaléniyu}, once more
  \item \emph{Say it:} say \emph{k sozhaléniyu}
  \item \emph{Recall:} say the Russian for probably, then the Russian for maybe, then say \emph{k sozhaléniyu} again
\end{itemize}

\begin{tcolorbox}[breakable,colback=teal!4,colframe=teal!35!black,title={Before you move on}]
What does \ru{к} \ru{сожалению} mean? (\textquotedblleft{}Unfortunately\textquotedblright{}.) Say it once more.
\end{tcolorbox}
\section[\texorpdfstring{\ru{к} \ru{счастью} (k schástyu)}{к счастью (k schástyu)}]{\ru{к} \ru{счастью} (k schástyu) — fortunately}

\label{lesson:RU-C215-k-schastyu}



\begin{quote}
Before the new one: say the Russian for maybe, then the Russian for unfortunately.
\end{quote}

\subsection*{You'll want to know: \ru{к} \ru{счастью}}
\textbf{\ru{к} \ru{счастью}} (\emph{k schástyu}) — \textquotedblleft{}fortunately\textquotedblright{}.

\begin{grammarlens}[title={What you've built}]
One more word for this chapter.
\end{grammarlens}

\subsection*{Your turn}
\begin{itemize}
\raggedright
  \item \emph{Say it:} \emph{k schástyu}
  \item \emph{Say it:} \emph{k schástyu}, once more
  \item \emph{Say it:} say \emph{k schástyu}
  \item \emph{Recall:} say the Russian for maybe, then the Russian for unfortunately, then say \emph{k schástyu} again
\end{itemize}

\begin{tcolorbox}[breakable,colback=teal!4,colframe=teal!35!black,title={Before you move on}]
What does \ru{к} \ru{счастью} mean? (\textquotedblleft{}Fortunately\textquotedblright{}.) Say it once more.
\end{tcolorbox}
